\scanpage{37}
\mspage{36}
% Reused typed sheet, read after rotation through 180 degrees.
\begin{center}-3.9-\end{center}
leaves $x$ fixed. Secondly, $3)\Rightarrow1)$ is an immediate consequence of Proposition 2.1.

Thirdly, we will prove that $2)\Rightarrow3)$. By virtue of 2), all the points $x_i$ cannot lie in one hyperplane, hence we can choose $n+1$ of the $x_i$ which are not in one hyperplane, say $x^{(0)},x^{(1)},\ldots,x^{(n)}$. Without loss of generality, we may assume that $x^{(i)}$ has homogeneous coordinates $x_j^{(i)}=\delta_{ij}$. Now for each $n+k$ between $n+1$ and $m$, let $S'_k$ be the set of integers $i$ such that
\[
D_{0,1,\ldots,\hat i,\ldots,n,n+k}\ne0,
\]
i.e. $x^{(n+k)}$ is not in the hyperplane spanned by $x^{(0)},\ldots,\widehat{x^{(i)}},\ldots,x^{(n)}$. Then I claim that there is no partition of the set $\{0,1,\ldots,n\}$ into 2 disjoint subsets $T'$ and $T''$ such that every $S'_k$ is contained in $T'$ or $T''$. For if there were, and if one let $L'$ (resp.$L''$) be the linear subspace defined by $X_i=0$ for all $i\in T'$ (resp. $i\in T''$), then every point $x^{(k)}$ would be in $L'\cup L''$. It follows immediately that a suitable set of subsets $S_i\subset S'_i$ is an R-partition $R$ and that $x\in\mathcal U_R$. \uline{QED}

From this Proposition, it is clear that $\mathcal U_{\mathrm{reg}}$ is precisely the set of pre-stable points of $(\mathbb P_n)^{m+1}$, in the sense of Chapter 1.

Our next step is to construct a geometric quotient of $\mathcal U_{\mathrm{reg}}$ by $\mathrm{PGL}(n)$: let $\mathcal U_1,\ldots,\mathcal U_N$ be the open subsets $\mathcal U_R$ of $\mathcal U_{\mathrm{reg}}$ and the subsets obtained from these by permuting the coordinates. Let $(Z_i,\phi_i)$ be the geometric quotient of $\mathcal U_i$ by $\mathrm{PGL}(n)$. For all pairs $i,j$, $\mathcal U_i\cap\mathcal U_j$ is an invariant open subset in $\mathcal U_i$ and in $\mathcal U_j$.
